\documentclass[11pt]{article}
\usepackage[margin=1.1in]{geometry}
\usepackage{amsmath,amssymb,amsthm}
\usepackage{booktabs}
\usepackage[hidelinks]{hyperref}

\newtheorem{theorem}{Theorem}
\newtheorem{lemma}[theorem]{Lemma}
\newtheorem{proposition}[theorem]{Proposition}
\newtheorem{corollary}[theorem]{Corollary}
\theoremstyle{remark}
\newtheorem{remark}[theorem]{Remark}

\newcommand{\ZZ}{\mathbb{Z}}
\newcommand{\CC}{\mathbb{C}}
\newcommand{\ip}[2]{\langle #1 | #2 \rangle}

\title{No MUB quadruple in dimension six from Butson Hadamard
matrices of root order~24\\[1ex]
\large A certified finite reduction (draft for verification)}
\author{D.~L.~Yonge-Mallo}
\date{2026-09-10 (draft v0.1)}

\begin{document}
\maketitle

\begin{abstract}
We prove that no four mutually unbiased bases of $\CC^6$ have the form
$\{I, H_1/\sqrt6, H_2/\sqrt6, H_3/\sqrt6\}$ with every entry of $H_1,
H_2, H_3$ a $24$th root of unity; equivalently, no MU quadruple in
dimension six can be built from Butson Hadamard matrices in
$\mathrm{BH}(6,24)$. The proof is a reduction of the statement to a
finite assertion (Proposition~\ref{prop:reduction}), checked by an
exhaustive exact computation in $\ZZ[\zeta_{24}]$ carried out by three
independently written programs. This confirms, in certified exact
arithmetic, the numerical search of Bengtsson et
al.~\cite{bengtsson2007} at $24$th roots, and extends the published
theorem record, which stops at $\mathrm{BH}(6,12)$
\cite[Thm.~7.15]{mcnulty2026}. As a by-product we obtain the first
classification of $\mathrm{BH}(6,24)$: it has exactly $25$
monomial-equivalence classes ($14$ up to transpose). This note contains
the complete reduction argument and the algorithm contract needed to
re-implement the computation independently; the accompanying repository
contains code, deterministic logs, and crosswalks to independently
published data.
\end{abstract}

\section{Introduction}

Two orthonormal bases $\{e_i\}$, $\{f_j\}$ of $\CC^d$ are
\emph{mutually unbiased} (MU) if $|\ip{e_i}{f_j}|^2 = 1/d$ for all
$i,j$. In every prime-power dimension exactly $d+1$ MU bases exist;
dimension six is the first open case, and Zauner
\cite{zauner1999} conjectured that no more than three MU bases exist in
$\CC^6$. Writing bases in Hadamard form, a hypothetical fourth basis
yields an \emph{MU quadruple} $\{I, H_1/\sqrt6, H_2/\sqrt6,
H_3/\sqrt6\}$ where each $H_k$ is a complex Hadamard matrix of order
six and the three are pairwise mutually unbiased. Progress on the
conjecture typically takes the form of excluding strata of the complex
Hadamard landscape from participating in a quadruple.

A \emph{Butson Hadamard matrix} $H \in \mathrm{BH}(d, N)$ is a $d
\times d$ matrix with all entries $N$th roots of unity and $H H^* = d
I$ \cite{butson1962}. These discrete strata admit finite exact
exhaustion. Bengtsson et al.~\cite{bengtsson2007} searched them
numerically in 2007: they enumerated the order-six Hadamard matrices
with $12$th- and $24$th-root entries, the vectors unbiased to each, the
third bases these admit, and whether any two third bases yield four MU
bases, finding only triplets. The survey \cite{mcnulty2026} records the
resulting theorem only at root order $12$ (its Theorem~7.15). The
present note upgrades the $24$th-root exclusion to a certified theorem.

\begin{theorem}\label{thm:main}
There is no MU quadruple $\{I, H_1/\sqrt6, H_2/\sqrt6, H_3/\sqrt6\}$ in
$\CC^6$ with $H_1, H_2, H_3 \in \mathrm{BH}(6, 24)$.
\end{theorem}

Since $\zeta_m = \zeta_{24}^{24/m}$ embeds the $m$th roots into the
$24$th roots for every $m \mid 24$, the theorem covers all entries
drawn from roots of unity of order dividing $24$, in any mixture
(Corollary~\ref{cor:divisors}), and a small additional computation
extends the exclusion to every single root order $N \le 12$
(Corollary~\ref{cor:small}). The theorem says nothing about other root
orders or about the continuous Hadamard families, and so does not
resolve Zauner's conjecture. The stratum is not vacuous: it contains
genuine MU triplets, and the computation in fact enumerates all
$1{,}920$ of them and shows none extends.

\section{Notation and elementary reductions}

Fix $N = 24$, $\omega = e^{2\pi i/24}$, and $d = 6$. We work with
\emph{exponent vectors} $a \in (\ZZ/24)^6$, representing the unimodular
vector $\omega^a = (\omega^{a_0}, \dots, \omega^{a_5}) \in \CC^6$. A
vector (or matrix row or column) is \emph{dephased} if $a_0 = 0$. For
$a \in (\ZZ/24)^6$ put
\[
  S(a) \;=\; \sum_{j=0}^{5} \omega^{a_j} \;\in\; \ZZ[\omega].
\]
We say $a$ satisfies the \emph{$Z$-condition} if $S(a) = 0$ and the
\emph{$M$-condition} if $|S(a)|^2 = 6$. Both conditions are invariant
under adding a constant to all coordinates of $a$, since this
multiplies $S(a)$ by a root of unity. For $u, w \in (\ZZ/24)^6$,
\[
  \ip{\omega^u}{\omega^w} \;=\; \sum_j \overline{\omega^{u_j}}\,
  \omega^{w_j} \;=\; S(w - u),
\]
so two unimodular root-valued vectors are \emph{orthogonal} iff their
exponent difference satisfies the $Z$-condition, and \emph{unbiased}
(for bases scaled by $1/\sqrt6$) iff it satisfies the $M$-condition.
All tests are exact: $\ZZ[\omega] \cong \ZZ[x]/\Phi_{24}(x)$ with
$\Phi_{24}(x) = x^8 - x^4 + 1$, so $S(a)$ and $|S(a)|^2 = S(a)
\overline{S(a)}$ are integer vectors of length $8$, and the conditions
are integer identities. Define the dephased ray classes
\[
  Z \;=\; \{a \in \{0\} \times (\ZZ/24)^5 : S(a) = 0\}, \qquad
  M \;=\; \{a \in \{0\} \times (\ZZ/24)^5 : |S(a)|^2 = 6\}.
\]

Note that any unimodular vector is automatically unbiased to every
standard basis vector, so unbiasedness to $I$ is free, and that a
$6\times 6$ matrix with unimodular entries has orthogonal columns iff
it has orthogonal rows (both say $\tfrac{1}{\sqrt6}H$ is unitary).

\begin{lemma}[Normal form]\label{lem:normal}
If an MU quadruple as in Theorem~\ref{thm:main} exists, then one exists
in which additionally $H_1$ is dephased (first row and first column all
ones) and every column of $H_2$ and $H_3$ is dephased.
\end{lemma}

\begin{proof}
Let $D_L$ be the diagonal unitary whose $j$th entry is the conjugate of
the $j$th entry of the first column of $H_1$; its entries are $24$th
roots. Replacing each basis $B$ of the quadruple by $D_L B$ preserves
all MU relations (a common unitary) and maps the standard basis to
itself up to phases; it multiplies each $H_k$ on the left by $D_L$,
keeping all entries in $\mu_{24}$ and making the first column of $H_1$
all ones. Next, multiplying any $H_k$ on the \emph{right} by a diagonal
matrix of $24$th roots rephases its columns, which changes neither the
basis it represents nor membership in $\mathrm{BH}(6,24)$: use this to
make the first row of $H_1$ all ones, and to dephase every column of
$H_2$ and $H_3$.
\end{proof}

\begin{lemma}[First basis]\label{lem:first}
The map sending a dephased $H \in \mathrm{BH}(6,24)$ to the set of its
five non-initial rows is a bijection between dephased
$\mathrm{BH}(6,24)$ matrices with sorted rows and $5$-element subsets
$R \subseteq Z$ whose pairwise differences satisfy the $Z$-condition.
Permuting the non-initial rows of $H_1$ preserves the quadruple
property, so it suffices to enumerate the subsets $R$.
\end{lemma}

\begin{proof}
A dephased $H$ has first row $\omega^0$ (all ones) and every row
$r$ with $r_0 = 0$. Orthogonality of row $r$ against the all-ones row
is $S(r) = 0$, i.e., $r \in Z$; orthogonality of rows $r, r'$ is the
$Z$-condition on $r' - r$, which is again dephased. Conversely such a
subset, together with the all-ones row, has pairwise orthogonal
unimodular rows, hence lies in $\mathrm{BH}(6,24)$. Row permutations
fixing the first row amount to left multiplication by a permutation
matrix, which fixes the standard basis and every MU relation.
\end{proof}

\begin{lemma}[Second and third bases]\label{lem:completion}
Fix a dephased $H_1$ with columns $h^{(0)} = 0, h^{(1)}, \dots,
h^{(5)}$ (exponent vectors), and define the \emph{completion rays}
\[
  U(H_1) \;=\; \bigl\{c \in M :\ c - h^{(i)} \text{ satisfies the
  $M$-condition for } i = 1, \dots, 5 \bigr\}.
\]
Then the bases unbiased to both $I$ and $H_1/\sqrt6$ that are
representable with columns in $\mathrm{BH}(6,24)$ form (after column
dephasing and up to column order) exactly the $6$-element subsets $C
\subseteq U(H_1)$ whose pairwise differences satisfy the
$Z$-condition.
\end{lemma}

\begin{proof}
By Lemma~\ref{lem:normal} each column $c$ of such a basis matrix may be
taken dephased with entries in $\mu_{24}$. Unbiasedness against the
column $h^{(0)} = 0$ of $H_1$ is the $M$-condition on $c - 0$, i.e., $c
\in M$; against the remaining columns it is the $M$-condition on $c -
h^{(i)}$ (the $M$-condition is phase-invariant, so the non-dephased
difference is immaterial). Hence each column lies in $U(H_1)$.
Orthogonality of two columns is the $Z$-condition on their difference.
Distinct columns have distinct dephased exponent vectors (equal ones
are not orthogonal), so the six columns form a $6$-subset; column order
is immaterial for the basis. The converse direction is immediate.
\end{proof}

\begin{lemma}[Pair test]\label{lem:pair}
With $H_1$ fixed as above, two bases given by $6$-cliques $C_2 \ne C_3
\subseteq U(H_1)$ are mutually unbiased iff every difference $c' - c$
with $c \in C_2$, $c' \in C_3$ satisfies the $M$-condition.
\end{lemma}

\begin{proof}
Immediate from the correspondence between unbiasedness and the
$M$-condition on exponent differences. (Note $C_2 = C_3$ can never
yield an MU pair: a basis is not unbiased to itself, and indeed $S(0) =
6$ fails the $M$-condition.)
\end{proof}

\begin{proposition}[Finite reduction]\label{prop:reduction}
Theorem~\ref{thm:main} holds if and only if for every $5$-subset $R$ as
in Lemma~\ref{lem:first} (equivalently, every dephased $H_1$), no two
distinct $6$-cliques $C_2, C_3 \subseteq U(H_1)$ have all $36$
cross-differences satisfying the $M$-condition.
\end{proposition}

\begin{proof}
Lemmas~\ref{lem:normal}--\ref{lem:pair} map any MU quadruple in the
stratum to such a configuration and conversely.
\end{proof}

\section{The computation}\label{sec:computation}

The finite assertion of Proposition~\ref{prop:reduction} is checked by
the following five-stage algorithm; the stage-by-stage counts are part
of the claim and serve as checkpoints for independent
re-implementations.

\begin{enumerate}
\item \textbf{Ray classification.} For each of the $24^5 = 7{,}962{,}624$
  dephased exponent vectors $a$, decide the $Z$- and $M$-conditions
  exactly in $\ZZ[x]/\Phi_{24}(x)$. Result: $|Z| = 8{,}350$ and $|M| =
  19{,}320$.
\item \textbf{Hadamard enumeration.} Enumerate all $5$-subsets of $Z$
  with pairwise differences in the $Z$-condition. Result: exactly
  $7{,}584$ dephased $\mathrm{BH}(6,24)$ row sets.
\item \textbf{Completion rays.} For each row set, compute $U(H_1)$ by
  the column tests of Lemma~\ref{lem:completion}. Census over the
  $7{,}584$ matrices, written $|U|/\#\text{bases}$:
  $0/0: 6{,}852$; \ $6/1: 240$; \ $12/4: 60$; \ $24/4: 360$; \
  $40/0: 72$.
\item \textbf{Completion bases.} Enumerate the $6$-cliques of
  Lemma~\ref{lem:completion}. Result: $1{,}920$ bases, spread over
  $660$ matrices ($240 \times 1 + 60 \times 4 + 360 \times 4$),
  giving $420 \times \binom{4}{2} = 2{,}520$ candidate pairs.
\item \textbf{Pair refutation.} For each of the $2{,}520$ pairs,
  exhibit a cross pair $(c, c')$ whose difference fails the
  $M$-condition. All $2{,}520$ pairs are refuted; no MU quadruple
  exists in the stratum.
\end{enumerate}

Three independently written programs (available in the accompanying
repository, with deterministic logs and SHA-256 manifests) implement
this contract: the original search, an adversarial referee harness that
re-derives the $\Phi_{24}$ reduction table from a recurrence and
verifies all bases and refutations in a second exact representation,
and a third, code-independent enumerator. All agree on every count
above. For re-implementers preferring floating point: over all $24^5$
rays the smallest margin of a non-zero-sum ray from $0$ is $0.035$ and
of a non-unbiased ray from unbiasedness is $1.2 \times 10^{-3}$, with
no ray in an ambiguous band at double precision; the certified programs
use exact arithmetic throughout.

\section{Corollaries}

\begin{corollary}\label{cor:divisors}
No MU quadruple $\{I, H_1/\sqrt6, H_2/\sqrt6, H_3/\sqrt6\}$ exists in
which every entry of $H_1, H_2, H_3$ is a root of unity of order
dividing $24$ (the orders may vary from entry to entry).
\end{corollary}

\begin{proof}
All entries then lie in $\mu_{24}$, so each $H_k \in \mathrm{BH}(6,24)$
and Theorem~\ref{thm:main} applies.
\end{proof}

\begin{corollary}\label{cor:small}
For every single root order $N \le 12$, no MU quadruple exists with all
entries of $H_1, H_2, H_3$ being $N$th roots of unity.
\end{corollary}

\begin{proof}
For $N \in \{1,2,3,4,6,8,12\}$ this is Corollary~\ref{cor:divisors}.
For $N \in \{5, 7, 11\}$ a vanishing sum of six $N$th roots of unity
would have length divisible by the prime $N$ (Lam--Leung
\cite{lamleung2000}), so no row can be orthogonal to the all-ones row
and $\mathrm{BH}(6,N) = \emptyset$. For $N \in \{9, 10\}$, direct
computation (seconds; included in the repository) shows the analogue of
$M$ is empty: no dephased $N$th-root vector is unbiased to the all-ones
vector. Dephasing as in Lemma~\ref{lem:normal} stays inside $\mu_N$, so
no basis can be unbiased to both $I$ and a dephased $H_1$, and not even
an MU triple exists in these strata.
\end{proof}

\section{External anchors and the $\mathrm{BH}(6,24)$
classification}\label{sec:anchors}

Every stage of the computation reproduces independently published
data.

\emph{Hadamard enumeration.} Closing the $7{,}584$ dephased row sets
(and the analogous enumerations at smaller orders) under the full
residual group of monomial equivalence yields $1, 1, 4, 3, 11$
equivalence classes at $N = 3, 4, 6, 8, 12$, matching the published
classification of Lampio, \"Osterg\aa rd and Sz\"oll\H osi
\cite{lampio2020} at every order they tabulate, and $\mathbf{25}$
classes at $N = 24$ ($14$ up to transpose), for which no published
value exists. Of the $25$ classes, exactly $4$, $3$, and $11$ contain
representatives of $\mathrm{BH}(6,6)$, $\mathrm{BH}(6,8)$, and
$\mathrm{BH}(6,12)$ respectively, and $12$ are new at root order $24$.
Representatives with their completion data accompany this note.

\emph{Completion census.} Exactly five of the $25$ classes admit at
least one completion basis, with per-matrix base counts $4, 1, 1, 4,
4$; they are identified as the classes of $F(0,0)$, $F(1/6, 0)$,
$F^{T}(1/6, 0)$, $F(1/6, 1/12)$, and $D(1/8)$, matching Fig.~2 of
\cite{bengtsson2007} class for class (one further class has $40$
completion rays but no basis, consistently omitted there). The
agreement of the exact census with the independent 2007 numerical
census is a strong mutual consistency check.

\emph{Theorem record.} Restricting the pipeline to $N = 12$ reproduces
the setting of \cite[Thm.~7.15]{mcnulty2026}: $2{,}184$ dephased
matrices, completion census $0/0: 1{,}884$, $6/1: 240$, $12/4: 60$,
$480$ completion bases, and no unbiased pair.

\section{Verification}

The reduction argument of this note is elementary and short; we invite
readers to check it directly. The finite assertion can be verified
either by running the accompanying programs (minutes; deterministic
logs and hashes provided) or, preferably, by independent
re-implementation from the contract of Section~\ref{sec:computation},
using the stage counts and the anchors of Section~\ref{sec:anchors} as
checkpoints. The computations were authored and cross-verified by three
independent AI systems under a human-directed protocol described in the
repository; confirmations and refutations are equally welcome and will
be credited.

\begin{thebibliography}{9}

\bibitem{zauner1999} G.~Zauner, \emph{Quantendesigns: Grundz\"uge
einer nichtkommutativen Designtheorie}, PhD thesis, Universit\"at Wien,
1999.

\bibitem{butson1962} A.~T.~Butson, Generalized Hadamard matrices,
\emph{Proc. Amer. Math. Soc.} \textbf{13} (1962), 894--898.

\bibitem{bengtsson2007} I.~Bengtsson, W.~Bruzda, {\AA}.~Ericsson,
J.-{\AA}.~Larsson, W.~Tadej, K.~\.Zyczkowski, Mutually unbiased bases
and Hadamard matrices of order six, \emph{J. Math. Phys.} \textbf{48}
(2007), 052106; arXiv:quant-ph/0610161.

\bibitem{lamleung2000} T.~Y.~Lam, K.~H.~Leung, On vanishing sums of
roots of unity, \emph{J. Algebra} \textbf{224} (2000), 91--109.

\bibitem{lampio2020} P.~H.~J.~Lampio, P.~R.~J.~\"Osterg\aa rd,
F.~Sz\"oll\H osi, Orderly generation of Butson Hadamard matrices,
\emph{Math. Comp.} \textbf{89} (2020), 313--331; arXiv:1707.02287.

\bibitem{mcnulty2026} D.~McNulty, S.~Weigert, Mutually unbiased bases
in composite dimensions -- a review, \emph{Quantum} \textbf{10} (2026),
2051; arXiv:2410.23997.

\end{thebibliography}

\end{document}
